\documentclass[11pt,a4paper]{article}
\usepackage[utf8]{inputenc}
\usepackage[T1]{fontenc}
\usepackage[spanish,es-minimal,es-noshorthands]{babel}
\usepackage{lmodern}
\linespread{1.22}
\usepackage[a4paper,top=2.4cm,bottom=2.2cm,left=1.5cm,right=1.5cm,headheight=1.1cm,headsep=0.5cm,footskip=1cm]{geometry}
\usepackage{amsmath,amssymb,cancel}
\usepackage{icomma}
\usepackage{graphicx}
\usepackage[table]{xcolor}
\usepackage{tikz}\usetikzlibrary{calc}
\usepackage{pgfplots}\pgfplotsset{compat=1.17}
\usepackage{fontawesome5}
\usepackage{ccicons}
\usepackage{fancyhdr}
\usepackage{enumitem}
\usepackage{booktabs}
\usepackage{tabularx}
\usepackage{needspace}
\usepackage[most]{tcolorbox}
\usepackage{titlesec}
\usepackage[hidelinks]{hyperref}
\hypersetup{pdftitle={Apuntes de teoría. Unidad 3: Vectores y aplicaciones. Rectas, planos y haces de planos. 2.º de Bachillerato MII},pdfsubject={Matemáticas II},pdfkeywords={vectores en el espacio, dependencia lineal, bases, producto escalar, producto vectorial, producto mixto, ecuaciones de la recta, ecuaciones del plano, posiciones relativas, haces de planos}}

% colors de la web (FentMates, bloque de Geometría: naranja; es manté el nom fmgreen dels apunts de les altres unitats)
\definecolor{fmnavy}{HTML}{1B2748}
\definecolor{fmgreen}{HTML}{DC7A29}
\definecolor{fmgreenl}{HTML}{F79645}
\definecolor{fmyellow}{HTML}{FFC94D}
\definecolor{fmgold}{HTML}{E9A919}
\definecolor{fmaccent}{HTML}{FF7A50}
\definecolor{fmteal}{HTML}{2BB3A3}
\definecolor{fmdef}{HTML}{1D4F73}\definecolor{fmdefbg}{HTML}{EEF4F9}
\definecolor{fmprop}{HTML}{3B4252}\definecolor{fmpropbg}{HTML}{F3F4F7}
\definecolor{fmex}{HTML}{A3203A}\definecolor{fmexbg}{HTML}{FBF0F2}
\definecolor{fmmuted}{HTML}{5D6685}\definecolor{fmball0}{HTML}{E2483D}\definecolor{fmball1}{HTML}{2F6FD6}
\definecolor{fmsurface}{HTML}{EEF1F8}

\newcommand{\lt}{<}\newcommand{\gt}{>}
% fórmula en línia que no es parteix mai; si no cap en una línia, es posa centrada
\newcommand{\fmm}[1]{\setbox0=\hbox{$#1$}\ifdim\wd0>\linewidth\fmdisp{#1}\else\mbox{\box0}\fi}
% fórmula centrada; si és més ampla que la línia, s'encongeix fins que hi cap
\newcommand{\fmdisp}[1]{\[\setbox2=\hbox{$\displaystyle #1$}\ifdim\wd2>\linewidth\resizebox{\linewidth}{!}{\box2}\else\box2\fi\]}
\DeclareUnicodeCharacter{2212}{-}
\DeclareUnicodeCharacter{2225}{\ensuremath{\parallel}}\DeclareUnicodeCharacter{22A5}{\ensuremath{\perp}}\DeclareUnicodeCharacter{2226}{\ensuremath{\nparallel}}
\DeclareUnicodeCharacter{207A}{\ensuremath{^{+}}}\DeclareUnicodeCharacter{207B}{\ensuremath{^{-}}}
\setlength{\parindent}{0pt}
\setlength{\parskip}{12pt}
\setlength{\lineskip}{5pt}\setlength{\lineskiplimit}{5pt}  % més aire al voltant de les fraccions en línia
\setlength{\emergencystretch}{3em}
\setlist{itemsep=4pt,topsep=4pt}
\renewcommand{\thesection}{3.\arabic{section}}
\titleformat{\section}{\color{fmnavy}\Large\bfseries}{\colorbox{fmdef}{\color{white}\strut\ \thesection\ }}{0.6em}{}[{\color{fmdef}\titlerule[1.2pt]}]
\titlespacing*{\section}{0pt}{16pt}{8pt}
\titleformat*{\subsection}{\color{fmgreen}\large\bfseries}
\allowdisplaybreaks

% dibuix de la targeta de la unitat a la web (illustrations.js), en blanc
\newcommand{\fmcard}{\begin{tikzpicture}[x=0.062cm,y=-0.062cm,white,line width=3.2pt,line cap=round,line join=round]
  \begin{scope}[transparency group,opacity=0.15]
  \draw[line width=2.6pt,opacity=0.7] (110,70)--(110,8) (110,70)--(210,92) (110,70)--(40,108);
  \fill[opacity=0.14] (50,40)--(150,22)--(210,62)--(110,82)--cycle;
  \draw (50,40)--(150,22)--(210,62)--(110,82)--cycle;
  \draw (60,110)--(190,14);
  \fill (128,47) circle (5.5);
  \draw[line width=3.4pt] (128,47)--(128,18) (122,26)--(128,16)--(134,26);
  \node[anchor=base west,inner sep=0pt,font=\fontsize{22}{22}\selectfont] at (136,22) {$\vec n$};
  \node[anchor=base west,inner sep=0pt,font=\fontsize{22}{22}\selectfont] at (196,112) {$\pi$};
  \end{scope}
\end{tikzpicture}}
\newcommand{\fmbrand}{\raisebox{-0.25\height}{\includegraphics[height=0.5cm]{logo-fentmates.pdf}}}

% capçalera i peu
\pagestyle{fancy}\fancyhf{}
\fancyhead[L]{\small\fmbrand}
\fancyhead[R]{\small\color{fmmuted}2.º Bachillerato MII · Unidad 3: Vectores, rectas y planos}
\fancyfoot[L]{\scriptsize\color{fmmuted}\href{https://creativecommons.org/licenses/by-nc-sa/4.0/deed.ca}{\ccbyncsa}\ \ FentMates · fentmates.org · CC BY-NC-SA 4.0}
\fancyfoot[R]{\small\color{fmmuted}\thepage}
\renewcommand{\headrulewidth}{0.4pt}\renewcommand{\headrule}{\color{fmmuted!40}\hrule height \headrulewidth}
\fancypagestyle{primera}{\fancyhf{}\renewcommand{\headrulewidth}{0pt}
  \fancyfoot[L]{\scriptsize\color{fmmuted}\href{https://creativecommons.org/licenses/by-nc-sa/4.0/deed.ca}{\ccbyncsa}\ \ FentMates · fentmates.org · CC BY-NC-SA 4.0}
  \fancyfoot[R]{\small\color{fmmuted}\thepage}}

% caixes (mateixos colors que la web)
% quan una caixa es parteix entre dues pàgines, la vora del tall és discontínua
\tcbset{fmbase/.style={enhanced,breakable,arc=3mm,boxrule=1pt,left=3.5mm,right=3.5mm,top=2mm,bottom=2.5mm,before skip=10pt,after skip=10pt,fonttitle=\bfseries,
  frame style={fill=tcbcolframe},
  overlay first={\fill[tcbcolback] ([xshift=1pt]frame.south west) rectangle ([xshift=-1pt,yshift=1.2pt]frame.south east);
    \draw[tcbcolframe,line width=1pt,dash pattern=on 5pt off 3pt] (frame.south west)--(frame.south east);},
  overlay middle={\fill[tcbcolback] ([xshift=1pt]frame.south west) rectangle ([xshift=-1pt,yshift=1.2pt]frame.south east);
    \fill[tcbcolback] ([xshift=1pt]frame.north west) rectangle ([xshift=-1pt,yshift=-1.2pt]frame.north east);
    \draw[tcbcolframe,line width=1pt,dash pattern=on 5pt off 3pt] (frame.south west)--(frame.south east) (frame.north west)--(frame.north east);},
  overlay last={\fill[tcbcolback] ([xshift=1pt]frame.north west) rectangle ([xshift=-1pt,yshift=-1.2pt]frame.north east);
    \draw[tcbcolframe,line width=1pt,dash pattern=on 5pt off 3pt] (frame.north west)--(frame.north east);}}}
\newtcolorbox{fmdef}[1]{fmbase,colback=fmdefbg,colframe=fmdef,coltitle=white,title={#1}}
\newtcolorbox{fmprop}[1]{fmbase,colback=fmpropbg,colframe=fmprop,coltitle=white,title={#1},halign title=center}
\newtcolorbox{fmex}[1]{fmbase,before upper=\raggedright,colback=fmexbg,colframe=fmex,colbacktitle=fmexbg,coltitle=fmex,title={\faPencil*\ #1},
  titlerule=0.4pt,titlerule style={fmex!40,dashed}}
\newtcolorbox{fmwarn}[1]{fmbase,colback=fmsf!6,colframe=fmsf,coltitle=white,title={#1}}
\newtcolorbox{fmnote}{enhanced,breakable,colback=white,colframe=white,borderline west={3pt}{0pt}{fmaccent},arc=0pt,left=3mm,right=2mm,top=1mm,bottom=1mm,before skip=6pt,after skip=6pt,before upper=\raggedright}
\newcommand{\fmkey}[1]{\colorbox{fmyellow!45}{\strut #1}}
\newcommand{\fmuline}[1]{\textbf{\color{fmaccent}#1}}
\definecolor{fmok}{HTML}{15803D}\definecolor{fmev}{HTML}{2563EB}\definecolor{fmsf}{HTML}{C2410C}\definecolor{fmsi}{HTML}{7C3AED}\definecolor{fmas}{HTML}{A3203A}
\newcommand{\fmbadge}[2]{\tcbox[on line,colback=#1,colframe=#1,arc=2mm,boxsep=0pt,left=4pt,right=4pt,top=1.5pt,bottom=1.5pt]{\color{white}\footnotesize\bfseries #2}}
\newcommand{\fmchip}[1]{\tcbox[on line,colback=fmsf!10,colframe=fmsf!50,arc=2mm,boxsep=1pt,left=4pt,right=4pt]{#1}}
\newcommand{\fmlevel}[2]{\tcbox[on line,colback=white,colframe=fmmuted!40,arc=2mm,boxsep=2pt,halign=center]{\begin{tabular}{c}{\scriptsize\color{fmmuted}\MakeUppercase{#1}}\\#2\end{tabular}}}
% gràfiques
\pgfplotsset{fmaxis/.style={axis lines=middle,grid=major,grid style={fmmuted!15},axis line style={fmmuted,-latex},
  tick label style={font=\scriptsize,color=fmmuted},clip=true,unbounded coords=jump,restrict y to domain=-200:200,
  every axis plot/.append style={line join=round}},
  fmcurve/.style={fmaccent,line width=1.6pt,smooth}}
\tikzset{fmasy/.style={fmmuted,dashed,line width=0.9pt},fmguide/.style={fmaccent,dotted,line width=0.8pt},
  fmdot/.style={circle,fill=fmaccent,inner sep=0pt,minimum size=5pt},fmopen/.style={circle,draw=fmaccent,fill=white,line width=1pt,inner sep=0pt,minimum size=5pt}}
\newcommand{\fmpau}[1]{\tcbox[on line,colback=fmnavy,colframe=fmnavy,arc=2mm,boxsep=0pt,left=4pt,right=4pt,top=1.5pt,bottom=1.5pt]{\color{white}\scriptsize\bfseries\MakeUppercase{#1}}\ }

\begin{document}
\thispagestyle{primera}
\begin{tikzpicture}[remember picture,overlay]
  \fill[fmgreen] (current page.north west) rectangle ([yshift=-10.1cm]current page.north east);
  \begin{scope}
    \clip (current page.north west) rectangle ([yshift=-10.1cm]current page.north east);
    \node[anchor=north west,inner sep=0pt] at ([xshift=6.6cm,yshift=-1.3cm]current page.north west) {\fmcard};
  \end{scope}
  \node[anchor=north west,inner sep=0pt] at ([xshift=1.5cm,yshift=-1.1cm]current page.north west) {\tcbox[colback=white,colframe=white,arc=2.5mm,boxsep=0pt,left=7pt,right=7pt,top=5pt,bottom=5pt]{\includegraphics[height=0.75cm]{logo-fentmates.pdf}}};
\end{tikzpicture}
\vspace*{0.0cm}
\tcbox[colback=white!15!fmgreen,colframe=white!15!fmgreen,arc=2.5mm,boxsep=0pt,left=6pt,right=6pt,top=3pt,bottom=3pt]{\color{white}\small\bfseries Apuntes de teoría · 2.º Bachillerato · Matemáticas II}

\vspace{4pt}
{\raggedright\color{white}\fontsize{25}{30}\selectfont\bfseries{\color{fmyellow}Unidad 3} · Vectores y aplicaciones. Rectas, planos y haces de planos\par}
\vspace{6pt}
\begin{minipage}{\linewidth}\color{white!88!fmgreen}\small Teoría y ejemplos resueltos paso a paso: vectores en el espacio y sus aplicaciones (punto medio, simétrico, puntos alineados, baricentro), dependencia lineal y bases, los productos escalar, vectorial y mixto (ángulos, áreas y volúmenes), todas las ecuaciones de la recta y del plano, las posiciones relativas de rectas y planos, la discusión de tres planos con el teorema de Rouché y los haces de planos.\end{minipage}\par
% l'índex comença sempre a la mateixa altura, per davall de la franja (com en les unitats 5-10)
\vspace*{\dimexpr 8.6cm-\pagetotal\relax}

\begin{tcolorbox}[colback=fmsurface,colframe=fmsurface,arc=3mm,left=4mm,right=4mm,top=2mm,bottom=2mm]
\small\textbf{Índice}\par\smallskip
\begin{minipage}[t]{0.49\linewidth}\small \mbox{\hyperlink{sec1}{\textbf{\color{fmgreen}3.1}\ \ Vectores en el espacio}}\\[3pt]
\mbox{\hyperlink{sec2}{\textbf{\color{fmgreen}3.2}\ \ Aplicaciones de los vectores}}\\[3pt]
\mbox{\hyperlink{sec3}{\textbf{\color{fmgreen}3.3}\ \ Dependencia lineal y bases}}\\[3pt]
\mbox{\hyperlink{sec4}{\textbf{\color{fmgreen}3.4}\ \ Producto escalar}}\\[3pt]
\mbox{\hyperlink{sec5}{\textbf{\color{fmgreen}3.5}\ \ Producto vectorial}}\\[3pt]
\mbox{\hyperlink{sec6}{\textbf{\color{fmgreen}3.6}\ \ Producto mixto}}\\[3pt]\end{minipage}\hfill\begin{minipage}[t]{0.49\linewidth}\small \mbox{\hyperlink{sec7}{\textbf{\color{fmgreen}3.7}\ \ Ecuaciones de la recta}}\\[3pt]
\mbox{\hyperlink{sec8}{\textbf{\color{fmgreen}3.8}\ \ Ecuaciones del plano}}\\[3pt]
\hyperlink{sec9}{\textbf{\color{fmgreen}3.9}\ \ Posiciones relativas de rectas y planos}\\[3pt]
\mbox{\hyperlink{sec10}{\textbf{\color{fmgreen}3.10}\ \ Posición relativa de tres planos}}\\[3pt]
\mbox{\hyperlink{sec11}{\textbf{\color{fmgreen}3.11}\ \ Haces de planos}}\\[3pt]\end{minipage}
\end{tcolorbox}

\Needspace{18\baselineskip}
\section{Vectores en el espacio}\hypertarget{sec1}{}

\Needspace{7\baselineskip}
\begin{fmdef}{Definiciones}
\begin{itemize}
\item Un \textbf{sistema de referencia} del espacio está formado por un punto \fmm{O} (el \textbf{origen}) y tres vectores \fmm{\vec i=(1,0,0)}, \fmm{\vec j=(0,1,0)}, \fmm{\vec k=(0,0,1)}, que determinan los ejes \fmm{OX}, \fmm{OY} y \fmm{OZ}. Cada punto \fmm{P} tiene tres coordenadas \fmm{P(x,y,z)}, que son las componentes de su \textbf{vector de posición} \fmm{\overrightarrow{OP}=x\,\vec i+y\,\vec j+z\,\vec k}.
\item Las \textbf{componentes} del vector \fmm{\overrightarrow{AB}} son las coordenadas del extremo menos las del origen:

\fmdisp{A(a_1,a_2,a_3),\ B(b_1,b_2,b_3)\ \Rightarrow\ \overrightarrow{AB}=(b_1-a_1,\ b_2-a_2,\ b_3-a_3)}

\item El \textbf{módulo} es la longitud del vector:

\fmdisp{|\vec v|=\sqrt{v_1^2+v_2^2+v_3^2}\qquad\qquad d(A,B)=|\overrightarrow{AB}|=\sqrt{(b_1-a_1)^2+(b_2-a_2)^2+(b_3-a_3)^2}}

\item Como en el plano, dos vectores son \textbf{iguales} (equipolentes) si tienen el mismo módulo, dirección y sentido, es decir, las mismas componentes. Un vector es \textbf{unitario} si tiene módulo 1; para \textbf{normalizar} un vector lo dividimos por su módulo: \fmm{\dfrac{\vec v}{|\vec v|}}.
\end{itemize}
\end{fmdef}

\Needspace{7\baselineskip}
\begin{fmprop}{Operaciones con vectores}
Si \fmm{\vec u=(u_1,u_2,u_3)}, \fmm{\vec v=(v_1,v_2,v_3)} y \fmm{k\in\mathbb R}, las operaciones se hacen componente a componente:

\fmdisp{\vec u+\vec v=(u_1+v_1,\ u_2+v_2,\ u_3+v_3)\qquad\qquad k\,\vec u=(k\,u_1,\ k\,u_2,\ k\,u_3)}

Gráficamente, \fmm{\vec u+\vec v} es la diagonal del paralelogramo que forman \fmm{\vec u} y \fmm{\vec v}; \fmm{k\,\vec u} tiene la misma dirección que \fmm{\vec u}, el mismo sentido si \fmm{k\gt0} y sentido contrario si \fmm{k\lt0}. Las propiedades (conmutativa, asociativa, distributivas\dots{}) son las mismas que en el plano.
\end{fmprop}

\Needspace{7\baselineskip}
\begin{fmex}{Ejemplo resuelto}
\begin{enumerate}[label=\textbf{\arabic*.},leftmargin=*]
\item Sean \fmm{A(2,-1,3)} y \fmm{B(5,1,-3)}. Calcula \fmm{\overrightarrow{AB}}, la distancia entre \fmm{A} y \fmm{B} y un vector unitario con la misma dirección y sentido que \fmm{\overrightarrow{AB}}.

\fmm{\overrightarrow{AB}=\big(5-2,\ 1-(-1),\ -3-3\big)=(3,\ 2,\ -6)}

\fmm{d(A,B)=|\overrightarrow{AB}|=\sqrt{3^2+2^2+(-6)^2}=\sqrt{9+4+36}=\sqrt{49}=7}

\fmm{\vec w=\dfrac{\overrightarrow{AB}}{|\overrightarrow{AB}|}=\left(\dfrac37,\ \dfrac27,\ -\dfrac67\right)}. Comprobación: \fmm{\dfrac{9+4+36}{49}=1\ \checkmark}
\item Si \fmm{\vec u=(1,-2,3)} y \fmm{\vec v=(2,0,-1)}, calcula \fmm{2\vec u-3\vec v}.

\fmm{2\vec u-3\vec v=(2,-4,6)-(6,0,-3)=(-4,\ -4,\ 9)}
\end{enumerate}
\end{fmex}

\Needspace{5\baselineskip}
\section{Aplicaciones de los vectores}\hypertarget{sec2}{}

\Needspace{7\baselineskip}
\begin{fmdef}{Fórmulas}
\begin{itemize}
\item \textbf{Punto medio} del segmento \fmm{AB}:
          \fmm{\ M=\left(\dfrac{a_1+b_1}{2},\ \dfrac{a_2+b_2}{2},\ \dfrac{a_3+b_3}{2}\right)}
\item \textbf{Simétrico} de \fmm{A} respecto de un punto \fmm{M}: como \fmm{M} es el punto medio de \fmm{A} y \fmm{A'},
          \fmm{\ A'=2M-A=(2m_1-a_1,\ 2m_2-a_2,\ 2m_3-a_3)}
\item \textbf{Dividir un segmento en \fmm{n} partes iguales}: los puntos de división son
          \fmm{\ P_k=A+\dfrac{k}{n}\,\overrightarrow{AB}}, \fmm{\ k=1,\dots,n-1}.
\item Tres puntos \fmm{A}, \fmm{B}, \fmm{C} están \textbf{alineados} si \fmm{\overrightarrow{AB}} y \fmm{\overrightarrow{AC}} son \textbf{proporcionales} (tienen la misma dirección):\\[2ex]
          \fmm{\ \dfrac{c_1-a_1}{b_1-a_1}=\dfrac{c_2-a_2}{b_2-a_2}=\dfrac{c_3-a_3}{b_3-a_3}} (equivalentemente, \fmm{\mathrm{rg}\big(\overrightarrow{AB},\overrightarrow{AC}\big)=1}).
          \newpage
\item \textbf{Baricentro} de un triángulo \fmm{ABC} (punto de corte de las medianas):
          \fmm{\ G=\left(\dfrac{a_1+b_1+c_1}{3},\ \dfrac{a_2+b_2+c_2}{3},\ \dfrac{a_3+b_3+c_3}{3}\right)}
\item \textbf{Paralelogramo} \fmm{ABCD}: \fmm{\overrightarrow{AB}=\overrightarrow{DC}}, y por tanto el cuarto vértice es \fmm{D=A+C-B}.
\end{itemize}
\end{fmdef}

\Needspace{7\baselineskip}
\begin{fmex}{Ejemplo resuelto}
\begin{enumerate}[label=\textbf{\arabic*.},leftmargin=*]
\item Calcula el punto medio de \fmm{A(2,-1,3)} y \fmm{B(4,5,-1)}, y el simétrico de \fmm{P(1,2,3)} respecto de \fmm{Q(2,0,-1)}.

\fmm{M=\left(\dfrac{2+4}{2},\ \dfrac{-1+5}{2},\ \dfrac{3+(-1)}{2}\right)=(3,\ 2,\ 1)}

\fmm{P'=2Q-P=(4-1,\ 0-2,\ -2-3)=(3,\ -2,\ -5)}. Comprobación: el punto medio de \fmm{P} y \fmm{P'} es \fmm{(2,0,-1)=Q\ \checkmark}
\item ¿Están alineados los puntos \fmm{A(1,2,-1)}, \fmm{B(3,1,1)} y \fmm{C(7,-1,5)}?

\fmm{\overrightarrow{AB}=(2,-1,2)}, \fmm{\quad\overrightarrow{AC}=(6,-3,6)}

\fmm{\dfrac62=\dfrac{-3}{-1}=\dfrac66=3\Rightarrow\overrightarrow{AC}=3\,\overrightarrow{AB}}. Son proporcionales: los puntos \textbf{están alineados}.
\item Divide el segmento de extremos \fmm{A(1,-2,4)} y \fmm{B(7,4,-2)} en tres partes iguales.

\fmm{\overrightarrow{AB}=(6,6,-6)\Rightarrow\dfrac13\overrightarrow{AB}=(2,2,-2)}

\fmm{P_1=A+\tfrac13\overrightarrow{AB}=(3,\ 0,\ 2)}, \fmm{\quad P_2=A+\tfrac23\overrightarrow{AB}=(5,\ 2,\ 0)}
\end{enumerate}
\end{fmex}

\Needspace{18\baselineskip}
\section{Dependencia lineal y bases}\hypertarget{sec3}{}

\Needspace{7\baselineskip}
\begin{fmdef}{Definiciones}
\begin{itemize}
\item Un vector \fmm{\vec w} es \textbf{combinación lineal} de \fmm{\vec u} y \fmm{\vec v} si existen números \fmm{a,b} tales que \fmm{\vec w=a\,\vec u+b\,\vec v}.
\item Unos vectores son \textbf{linealmente dependientes} si alguno es combinación lineal de los otros, y \textbf{linealmente independientes} en caso contrario. Geométricamente: dos vectores son dependientes si son \textbf{paralelos} (proporcionales); tres vectores son dependientes si son \textbf{coplanarios} (se pueden dibujar en un mismo plano).
\item El \textbf{rango} de un conjunto de vectores es el número máximo de vectores linealmente independientes que hay. Se calcula como el rango de la matriz que tiene los vectores por filas (unidad 1).
\item Una \textbf{base} del espacio es cualquier conjunto de \textbf{tres vectores linealmente independientes}. Entonces todo vector se escribe de manera única como combinación lineal de ellos: \fmm{\vec x=a\,\vec u+b\,\vec v+c\,\vec w}; los números \fmm{(a,b,c)} son las \textbf{coordenadas} de \fmm{\vec x} en esa base. La base \fmm{\{\vec i,\vec j,\vec k\}} es la \textbf{base canónica}.
\end{itemize}
\end{fmdef}

\Needspace{7\baselineskip}
\begin{fmprop}{Criterio del determinante}

\fmdisp{\{\vec u,\vec v,\vec w\}\ \text{es base}\iff\begin{vmatrix}u_1&u_2&u_3\\v_1&v_2&v_3\\w_1&w_2&w_3\end{vmatrix}\neq0}

Si el determinante es 0, los tres vectores son coplanarios (rango 2) o, incluso, paralelos (rango 1). En el espacio, \textbf{cuatro o más vectores siempre son linealmente dependientes}.
\end{fmprop}

\Needspace{7\baselineskip}
\begin{fmex}{Ejemplo resuelto}
\begin{enumerate}[label=\textbf{\arabic*.},leftmargin=*]
\item Estudia si \fmm{\vec u=(1,2,-1)}, \fmm{\vec v=(0,1,3)} y \fmm{\vec w=(2,5,1)} forman una base. Si no la forman, escribe \fmm{\vec w} como combinación lineal de \fmm{\vec u} y \fmm{\vec v}.

\fmm{\begin{vmatrix}1&2&-1\\0&1&3\\2&5&1\end{vmatrix}=1\cdot(1-15)-2\cdot(0-6)+(-1)\cdot(0-2)=-14+12+2=0}

El determinante es 0: \textbf{no forman una base} (son coplanarios). Como \fmm{\vec u} y \fmm{\vec v} no son proporcionales, el rango es 2.

\fmm{(2,5,1)=a(1,2,-1)+b(0,1,3)\Rightarrow\begin{cases}a=2\\2a+b=5\\-a+3b=1\end{cases}\Rightarrow a=2,\ b=1}. La tercera ecuación también se cumple: \fmm{-2+3=1\ \checkmark}

\fmm{\vec w=2\,\vec u+\vec v}
\item ¿Para qué valores de \fmm{m} forman una base los vectores \fmm{(1,1,m)}, \fmm{(m,1,1)} y \fmm{(1,m,1)}?

\fmm{\begin{vmatrix}1&1&m\\m&1&1\\1&m&1\end{vmatrix}=(1-m)-(m-1)+m(m^2-1)=(m-1)(m^2+m-2)=(m-1)^2(m+2)}

Se anula para \fmm{m=1} y \fmm{m=-2}. Por tanto, \textbf{forman base si \fmm{m\neq1} y \fmm{m\neq-2}}.
\end{enumerate}
\end{fmex}

\Needspace{7\baselineskip}
\section{Producto escalar}\hypertarget{sec4}{}

\Needspace{7\baselineskip}
\begin{fmdef}{Definición y expresión analítica}
El \textbf{producto escalar} de dos vectores es un \textbf{número}:

\fmdisp{\vec u\cdot\vec v=|\vec u|\,|\vec v|\cos\alpha\qquad\qquad \vec u\cdot\vec v=u_1v_1+u_2v_2+u_3v_3}

donde \fmm{\alpha} es el ángulo que forman (en una base ortonormal como la canónica, las dos expresiones coinciden).

\begin{itemize}
\item \textbf{Propiedades}: \fmm{\vec u\cdot\vec v=\vec v\cdot\vec u}; \fmm{\ \vec u\cdot(\vec v+\vec w)=\vec u\cdot\vec v+\vec u\cdot\vec w}; \fmm{\ (k\vec u)\cdot\vec v=k(\vec u\cdot\vec v)}; \fmm{\ \vec u\cdot\vec u=|\vec u|^2}.
\item \textbf{Ángulo entre dos vectores}: \fmm{\ \cos\alpha=\dfrac{\vec u\cdot\vec v}{|\vec u|\,|\vec v|}}, \fmm{\ 0^\circ\le\alpha\le180^\circ}.
\item \textbf{Perpendicularidad} (ortogonalidad): \fmm{\ \vec u\perp\vec v\iff\vec u\cdot\vec v=0}.
\item \textbf{Proyección} de \fmm{\vec u} sobre \fmm{\vec v}: la proyección escalar es \fmm{\dfrac{\vec u\cdot\vec v}{|\vec v|}} (con signo) y el vector proyección es \fmm{\dfrac{\vec u\cdot\vec v}{|\vec v|^2}\,\vec v}.
\item \textbf{Cosenos directores}: los cosenos de los ángulos que forma \fmm{\vec v} con los ejes son \fmm{\cos\alpha=\dfrac{v_1}{|\vec v|}}, \fmm{\cos\beta=\dfrac{v_2}{|\vec v|}}, \fmm{\cos\gamma=\dfrac{v_3}{|\vec v|}}, y cumplen \fmm{\cos^2\alpha+\cos^2\beta+\cos^2\gamma=1}.
\end{itemize}
\end{fmdef}

\Needspace{7\baselineskip}
\begin{fmex}{Ejemplo resuelto}
\begin{enumerate}[label=\textbf{\arabic*.},leftmargin=*]
\item Calcula el producto escalar y el ángulo de \fmm{\vec u=(2,-1,3)} y \fmm{\vec v=(1,4,-2)}.

\fmm{\vec u\cdot\vec v=2\cdot1+(-1)\cdot4+3\cdot(-2)=2-4-6=-8}

\fmm{|\vec u|=\sqrt{4+1+9}=\sqrt{14}}, \fmm{\quad|\vec v|=\sqrt{1+16+4}=\sqrt{21}}

\fmm{\cos\alpha=\dfrac{-8}{\sqrt{14}\sqrt{21}}=\dfrac{-8}{7\sqrt6}\approx-0{,}4666\Rightarrow\alpha\approx117^\circ48'43''}. Como el producto escalar es negativo, el ángulo es obtuso.
\item Halla \fmm{k} para que \fmm{\vec u=(k,2,-1)} y \fmm{\vec v=(3,k,4)} sean perpendiculares.

\fmm{\vec u\cdot\vec v=3k+2k-4=5k-4=0\Rightarrow k=\dfrac45}
\end{enumerate}
\end{fmex}

\Needspace{18\baselineskip}
\section{Producto vectorial}\hypertarget{sec5}{}

\Needspace{7\baselineskip}
\begin{fmdef}{Definición}
El \textbf{producto vectorial} \fmm{\vec u\times\vec v} es un \textbf{vector}:

\begin{itemize}
\item de \textbf{módulo} \fmm{|\vec u\times\vec v|=|\vec u|\,|\vec v|\sin\alpha};
\item de \textbf{dirección} perpendicular a \fmm{\vec u} y a \fmm{\vec v};
\item de \textbf{sentido} el de avance de un sacacorchos que gira de \fmm{\vec u} hacia \fmm{\vec v} (regla de la mano derecha).
\end{itemize}

Se calcula con un determinante «simbólico» desarrollado por la primera fila:

\fmdisp{\vec u\times\vec v=\begin{vmatrix}\vec i&\vec j&\vec k\\u_1&u_2&u_3\\v_1&v_2&v_3\end{vmatrix}=\left(\begin{vmatrix}u_2&u_3\\v_2&v_3\end{vmatrix},\ -\begin{vmatrix}u_1&u_3\\v_1&v_3\end{vmatrix},\ \begin{vmatrix}u_1&u_2\\v_1&v_2\end{vmatrix}\right)}

\end{fmdef}

\Needspace{7\baselineskip}
\begin{fmprop}{Propiedades y aplicaciones}
\begin{itemize}
\item \textbf{Anticonmutativa}: \fmm{\vec v\times\vec u=-(\vec u\times\vec v)}. \fmuline{¡El orden importa!}
\item \fmm{\vec u\times\vec v=\vec 0\iff\vec u\parallel\vec v} (en particular, \fmm{\vec u\times\vec u=\vec 0}). Es una manera rápida de ver si dos vectores son paralelos o si tres puntos están alineados.
\item Distributiva: \fmm{\vec u\times(\vec v+\vec w)=\vec u\times\vec v+\vec u\times\vec w}, \fmm{\ (k\vec u)\times\vec v=k(\vec u\times\vec v)}. \textbf{No} es asociativa.
\item \textbf{Vector perpendicular a dos vectores} dados: \fmm{\vec u\times\vec v}. Será la forma habitual de obtener el vector normal de un plano y el vector director de una recta dada por dos planos.
\item \textbf{Áreas}: \fmm{\ S_{\text{paralelogramo}}=|\overrightarrow{AB}\times\overrightarrow{AD}|}, \fmm{\quad S_{\text{triángulo}}=\dfrac12\,|\overrightarrow{AB}\times\overrightarrow{AC}|}.
\end{itemize}
\end{fmprop}

\Needspace{7\baselineskip}
\begin{fmex}{Ejemplo resuelto}
\begin{enumerate}[label=\textbf{\arabic*.},leftmargin=*]
\item Calcula \fmm{\vec u\times\vec v} si \fmm{\vec u=(1,2,-1)} y \fmm{\vec v=(3,0,2)}, y comprueba que es perpendicular a ambos.

\fmm{\vec u\times\vec v=\begin{vmatrix}\vec i&\vec j&\vec k\\1&2&-1\\3&0&2\end{vmatrix}=\left(\begin{vmatrix}2&-1\\0&2\end{vmatrix},\ -\begin{vmatrix}1&-1\\3&2\end{vmatrix},\ \begin{vmatrix}1&2\\3&0\end{vmatrix}\right)=(4,\ -5,\ -6)}

\fmm{(4,-5,-6)\cdot(1,2,-1)=4-10+6=0\ \checkmark\qquad(4,-5,-6)\cdot(3,0,2)=12+0-12=0\ \checkmark}
\item Calcula el área del triángulo de vértices \fmm{A(1,0,0)}, \fmm{B(0,2,0)} y \fmm{C(0,0,3)}.

\fmm{\overrightarrow{AB}=(-1,2,0)}, \fmm{\quad\overrightarrow{AC}=(-1,0,3)}

\fmm{\overrightarrow{AB}\times\overrightarrow{AC}=\begin{vmatrix}\vec i&\vec j&\vec k\\-1&2&0\\-1&0&3\end{vmatrix}=(6,\ 3,\ 2)}, \fmm{\quad|(6,3,2)|=\sqrt{36+9+4}=7}

\fmm{S=\dfrac12\cdot7=\dfrac72\ u^2}
\end{enumerate}
\end{fmex}

\Needspace{7\baselineskip}
\section{Producto mixto}\hypertarget{sec6}{}

\Needspace{7\baselineskip}
\begin{fmdef}{Definición y propiedades}
El \textbf{producto mixto} de tres vectores es el \textbf{número} \fmm{[\vec u,\vec v,\vec w]=\vec u\cdot(\vec v\times\vec w)}, que es igual al determinante que tiene los vectores por filas:

\fmdisp{[\vec u,\vec v,\vec w]=\begin{vmatrix}u_1&u_2&u_3\\v_1&v_2&v_3\\w_1&w_2&w_3\end{vmatrix}}
\newpage
\begin{itemize}
\item Por las propiedades de los determinantes: si permutamos dos vectores cambia de signo; \fmm{[\vec u,\vec v,\vec w]=[\vec v,\vec w,\vec u]=[\vec w,\vec u,\vec v]}.
\item \fmm{[\vec u,\vec v,\vec w]=0\iff} los tres vectores son \textbf{coplanarios} (linealmente dependientes).
\item Cuatro puntos \fmm{A,B,C,D} son \textbf{coplanarios} \fmm{\iff\big[\overrightarrow{AB},\overrightarrow{AC},\overrightarrow{AD}\big]=0}.
\end{itemize}
\end{fmdef}

\Needspace{7\baselineskip}
\begin{fmprop}{Volúmenes}

\fmdisp{V_{\text{paralelepípedo}}=\big|[\vec u,\vec v,\vec w]\big|\qquad\qquad V_{\text{tetraedro}\,ABCD}=\dfrac16\,\Big|\big[\overrightarrow{AB},\overrightarrow{AC},\overrightarrow{AD}\big]\Big|}

El valor absoluto es importante: el producto mixto puede ser negativo, pero un volumen no. El tetraedro es la sexta parte del paralelepípedo: el prisma triangular es la mitad del paralelepípedo y la pirámide, un tercio del prisma.
\end{fmprop}

\Needspace{7\baselineskip}
\begin{fmex}{Ejemplo resuelto}
\begin{enumerate}[label=\textbf{\arabic*.},leftmargin=*]
\item Calcula el volumen del paralelepípedo determinado por \fmm{\vec u=(1,0,2)}, \fmm{\vec v=(2,1,0)} y \fmm{\vec w=(0,3,1)}.

\fmm{[\vec u,\vec v,\vec w]=\begin{vmatrix}1&0&2\\2&1&0\\0&3&1\end{vmatrix}=1\cdot(1-0)-0+2\cdot(6-0)=13}

\fmm{V=|13|=13\ u^3}
\item Calcula el volumen del tetraedro de vértices \fmm{A(1,0,0)}, \fmm{B(0,1,0)}, \fmm{C(0,0,1)} y \fmm{D(2,2,2)}.

\fmm{\overrightarrow{AB}=(-1,1,0)}, \fmm{\ \overrightarrow{AC}=(-1,0,1)}, \fmm{\ \overrightarrow{AD}=(1,2,2)}

\fmm{\begin{vmatrix}-1&1&0\\-1&0&1\\1&2&2\end{vmatrix}=-1\cdot(0-2)-1\cdot(-2-1)+0=2+3=5}

\fmm{V=\dfrac{|5|}{6}=\dfrac56\ u^3}
\item Halla \fmm{k} para que los puntos \fmm{A(1,0,1)}, \fmm{B(2,1,0)}, \fmm{C(0,1,2)} y \fmm{D(k,2,1)} sean coplanarios.

\fmm{\overrightarrow{AB}=(1,1,-1)}, \fmm{\ \overrightarrow{AC}=(-1,1,1)}, \fmm{\ \overrightarrow{AD}=(k-1,2,0)}

\fmm{\begin{vmatrix}1&1&-1\\-1&1&1\\k-1&2&0\end{vmatrix}=1\cdot(0-2)-1\cdot\big(0-(k-1)\big)+(-1)\cdot\big(-2-(k-1)\big)=-2+k-1+k+1=2k-2}

\fmm{2k-2=0\Rightarrow k=1}
\end{enumerate}
\end{fmex}

\Needspace{18\baselineskip}
\section{Ecuaciones de la recta}\hypertarget{sec7}{}

\Needspace{7\baselineskip}
\begin{fmdef}{Una recta queda determinada por un punto y un vector director}
Sea \fmm{r} la recta que pasa por el punto \fmm{P(p_1,p_2,p_3)} con \textbf{vector director} \fmm{\vec d=(d_1,d_2,d_3)\neq\vec0}. Un punto \fmm{X(x,y,z)} está en \fmm{r} si \fmm{\overrightarrow{PX}} es proporcional a \fmm{\vec d}:

\begin{center}\renewcommand{\arraystretch}{1.5}\small\begin{tabularx}{\linewidth}{>{\raggedright\arraybackslash}p{4.2cm}>{\raggedright\arraybackslash}X}\toprule
\textbf{Ecuación} & \textbf{Expresión}\\\midrule
Vectorial & $(x,y,z)=(p_1,p_2,p_3)+\lambda\,(d_1,d_2,d_3),\ \lambda\in\mathbb R$\\[3pt]
Paramétricas & $\begin{cases}x=p_1+\lambda d_1\\y=p_2+\lambda d_2\\z=p_3+\lambda d_3\end{cases}$\\[3pt]
Continua & $\dfrac{x-p_1}{d_1}=\dfrac{y-p_2}{d_2}=\dfrac{z-p_3}{d_3}$\\[3pt]
Implícitas (o cartesianas) & $\begin{cases}A_1x+B_1y+C_1z+D_1=0\\A_2x+B_2y+C_2z+D_2=0\end{cases}$\\[3pt]
\bottomrule\end{tabularx}\end{center}

\begin{itemize}
\item Las \textbf{ecuaciones implícitas} describen la recta como la intersección de dos planos no paralelos. Se obtienen de la continua igualando las fracciones dos a dos (multiplicando en cruz).
\item \textbf{De las implícitas a las paramétricas}: resolvemos el sistema (compatible indeterminado con un parámetro), o bien tomamos \fmm{\vec d=\vec n_1\times\vec n_2} (producto vectorial de los vectores normales de los dos planos) y un punto cualquiera que cumpla las dos ecuaciones.
\item \textbf{Recta por dos puntos} \fmm{A} y \fmm{B}: punto \fmm{A} y vector director \fmm{\overrightarrow{AB}}.
\end{itemize}
\vspace*{0.2cm}
\begin{fmnote}\fmuline{Notación}: si una componente del vector director es 0, en la continua escribimos igualmente, por ejemplo, \fmm{\dfrac{x-1}{2}=\dfrac{y+3}{0}=\dfrac{z}{5}}. No es una división de verdad: indica que \fmm{y=-3} para todos los puntos de la recta.\end{fmnote}
\end{fmdef}

\Needspace{7\baselineskip}
\begin{fmex}{Ejemplo resuelto}
\begin{enumerate}[label=\textbf{\arabic*.},leftmargin=*]
\item Escribe todas las ecuaciones de la recta que pasa por \fmm{P(1,-2,3)} con vector director \fmm{\vec d=(2,1,-1)}.

Vectorial: \fmm{\ (x,y,z)=(1,-2,3)+\lambda(2,1,-1)}

Paramétricas: \fmm{\ \begin{cases}x=1+2\lambda\\y=-2+\lambda\\z=3-\lambda\end{cases}}\\

\vspace*{0.2cm}
Continua: \fmm{\ \dfrac{x-1}{2}=\dfrac{y+2}{1}=\dfrac{z-3}{-1}}

Implícitas: \fmm{\dfrac{x-1}{2}=\dfrac{y+2}{1}\Rightarrow x-1=2y+4\Rightarrow x-2y-5=0}; \fmm{\ \dfrac{y+2}{1}=\dfrac{z-3}{-1}\Rightarrow -y-2=z-3\Rightarrow y+z-1=0}.

\fmdisp{r:\begin{cases}x-2y-5=0\\y+z-1=0\end{cases}}

\item Escribe en forma paramétrica la recta \fmm{r:\begin{cases}x+y-z=1\\2x-y+z=2\end{cases}}

Vector director: \fmm{\ \vec n_1\times\vec n_2=\begin{vmatrix}\vec i&\vec j&\vec k\\1&1&-1\\2&-1&1\end{vmatrix}=(0,-3,-3)}, o más sencillo, \fmm{\vec d=(0,1,1)}.

Un punto: hacemos \fmm{z=0\Rightarrow\begin{cases}x+y=1\\2x-y=2\end{cases}\Rightarrow x=1,\ y=0}. Punto \fmm{P(1,0,0)}.

\fmm{r:\begin{cases}x=1\\y=\lambda\\z=\lambda\end{cases}}. (También lo obtenemos resolviendo el sistema con \fmm{z=\lambda}: sumando las ecuaciones, \fmm{3x=3}.)
\end{enumerate}
\end{fmex}

\Needspace{18\baselineskip}
\section{Ecuaciones del plano}\hypertarget{sec8}{}

\Needspace{7\baselineskip}
\begin{fmdef}{Un plano queda determinado por un punto y dos vectores directores no paralelos}
Sea \fmm{\pi} el plano que pasa por \fmm{P(p_1,p_2,p_3)} con vectores directores \fmm{\vec u} y \fmm{\vec v} (no paralelos). Un punto \fmm{X} está en \fmm{\pi} si \fmm{\overrightarrow{PX}} es combinación lineal de \fmm{\vec u} y \fmm{\vec v}:

\begin{center}\renewcommand{\arraystretch}{1.5}\small\begin{tabularx}{\linewidth}{>{\raggedright\arraybackslash}p{3.2cm}>{\raggedright\arraybackslash}X}\toprule
\textbf{Ecuación} & \textbf{Expresión}\\\midrule
Vectorial & $(x,y,z)=(p_1,p_2,p_3)+\lambda\,\vec u+\mu\,\vec v,\ \ \lambda,\mu\in\mathbb R$\\[3pt]
Paramétricas & $\begin{cases}x=p_1+\lambda u_1+\mu v_1\\y=p_2+\lambda u_2+\mu v_2\\z=p_3+\lambda u_3+\mu v_3\end{cases}$\\[3pt]
General (o implícita) & $\begin{vmatrix}x-p_1&y-p_2&z-p_3\\u_1&u_2&u_3\\v_1&v_2&v_3\end{vmatrix}=0\ \Longrightarrow\ Ax+By+Cz+D=0$\\[3pt]
Segmentaria & $\dfrac xa+\dfrac yb+\dfrac zc=1$ (corta los ejes en $(a,0,0)$, $(0,b,0)$, $(0,0,c)$)\\[3pt]
\bottomrule\end{tabularx}\end{center}

La ecuación general dice que \fmm{\overrightarrow{PX}}, \fmm{\vec u} y \fmm{\vec v} son coplanarios (producto mixto nulo).
\end{fmdef}

\Needspace{7\baselineskip}
\begin{fmprop}{El vector normal}
En la ecuación general \fmm{\pi:Ax+By+Cz+D=0}, el vector \fmm{\vec n=(A,B,C)} es \textbf{perpendicular} al plano: es el \textbf{vector normal}. Por tanto:

\begin{itemize}
\item \textbf{Plano por un punto con vector normal} \fmm{\vec n}: \fmm{\ A(x-p_1)+B(y-p_2)+C(z-p_3)=0}.
\item Si conocemos dos vectores directores, \fmm{\vec n=\vec u\times\vec v}.
\item De la general a las paramétricas: despejamos una incógnita y las otras dos son los parámetros (es un sistema de una ecuación con tres incógnitas).
\end{itemize}

\textbf{Casos típicos para determinar un plano}

\begin{center}\renewcommand{\arraystretch}{1.5}\small\begin{tabularx}{\linewidth}{>{\raggedright\arraybackslash}p{7.2cm}>{\raggedright\arraybackslash}X}\toprule
\textbf{Datos} & \textbf{Punto y vectores directores}\\\midrule
Tres puntos no alineados $A,B,C$ & $A;\ \overrightarrow{AB},\ \overrightarrow{AC}$\\[3pt]
Una recta $r(P,\vec d)$ y un punto $Q\notin r$ & $P;\ \vec d,\ \overrightarrow{PQ}$\\[3pt]
Dos rectas secantes $r(P,\vec d_r)$ y $s(Q,\vec d_s)$ & $P;\ \vec d_r,\ \vec d_s$\\[3pt]
Dos rectas paralelas $r(P,\vec d)$ y $s(Q,\vec d)$ & $P;\ \vec d,\ \overrightarrow{PQ}$\\[3pt]
Contiene $r(P,\vec d_r)$ y es paralelo a $s$ & $P;\ \vec d_r,\ \vec d_s$\\[3pt]
Pasa por $P$ y es perpendicular a la recta $r$ & $P;\ \vec n=\vec d_r$\\[3pt]
\bottomrule\end{tabularx}\end{center}
\end{fmprop}

\Needspace{7\baselineskip}
\begin{fmex}{Ejemplo resuelto}
\begin{enumerate}[label=\textbf{\arabic*.},leftmargin=*]
\item Halla la ecuación general del plano que pasa por \fmm{P(1,0,2)} con vectores directores \fmm{\vec u=(1,1,0)} y \fmm{\vec v=(0,2,-1)}.

\fmm{\begin{vmatrix}x-1&y&z-2\\1&1&0\\0&2&-1\end{vmatrix}=(x-1)(-1-0)-y(-1-0)+(z-2)(2-0)=0}

\fmm{-x+1+y+2z-4=0\Rightarrow x-y-2z+3=0}. Comprobación con \fmm{P}: \fmm{1-0-4+3=0\ \checkmark}
\item Halla el plano que pasa por \fmm{A(1,0,0)}, \fmm{B(0,2,0)} y \fmm{C(0,0,3)}.

Ya calculamos (apartado 3.5) que \fmm{\overrightarrow{AB}\times\overrightarrow{AC}=(6,3,2)}: es un vector normal.

\fmm{6(x-1)+3y+2z=0\Rightarrow 6x+3y+2z-6=0}

Como los puntos están en los ejes, también sale directamente de la segmentaria: \fmm{\dfrac x1+\dfrac y2+\dfrac z3=1}, y multiplicando por 6 obtenemos lo mismo.
\item Escribe en paramétricas el plano \fmm{x-2y+3z=4}.

Despejamos \fmm{x=4+2y-3z} y hacemos \fmm{y=\lambda}, \fmm{z=\mu}: \fmm{\ \begin{cases}x=4+2\lambda-3\mu\\y=\lambda\\z=\mu\end{cases}}
\end{enumerate}
\end{fmex}

\Needspace{18\baselineskip}
\section{Posiciones relativas de rectas y planos}\hypertarget{sec9}{}

\Needspace{16\baselineskip}
\subsection*{A. Dos rectas}

\Needspace{7\baselineskip}
\begin{fmprop}{Con los vectores directores y un vector que las une}
Sean \fmm{r} (punto \fmm{P}, director \fmm{\vec d_r}) y \fmm{s} (punto \fmm{Q}, director \fmm{\vec d_s}). Comparamos el rango de \fmm{\{\vec d_r,\vec d_s\}} con el de \fmm{\{\vec d_r,\vec d_s,\overrightarrow{PQ}\}}:

\begin{center}\renewcommand{\arraystretch}{1.5}\small\begin{tabularx}{\linewidth}{>{\raggedright\arraybackslash}p{3.2cm}>{\raggedright\arraybackslash}X>{\raggedright\arraybackslash}X>{\raggedright\arraybackslash}X}\toprule
\textbf{$\mathrm{rg}(\vec d_r,\vec d_s)$} & \textbf{$\mathrm{rg}(\vec d_r,\vec d_s,\overrightarrow{PQ})$} & \textbf{Posición} & \textbf{Rangos con las implícitas}\\\midrule
1 & 1 & \textbf{Coincidentes} & $\mathrm{rg}\,M=\mathrm{rg}\,M^*=2$\\[3pt]
1 & 2 & \textbf{Paralelas} & $\mathrm{rg}\,M=2,\ \mathrm{rg}\,M^*=3$\\[3pt]
2 & 2 & \textbf{Secantes} (se cortan en un punto) & $\mathrm{rg}\,M=\mathrm{rg}\,M^*=3$\\[3pt]
2 & 3 & \textbf{Se cruzan} (no están en un mismo plano) & $\mathrm{rg}\,M=3,\ \mathrm{rg}\,M^*=4$\\[3pt]
\bottomrule\end{tabularx}\end{center}

En la práctica: si \fmm{\vec d_r\parallel\vec d_s}, miramos si \fmm{P\in s} (coincidentes) o no (paralelas). Si no son paralelos, calculamos \fmm{\big[\vec d_r,\vec d_s,\overrightarrow{PQ}\big]}: si es 0 son \textbf{secantes} y si no, \textbf{se cruzan}. La última columna corresponde a las rectas dadas por implícitas: \fmm{M} es la matriz \fmm{4\times3} de los coeficientes de los cuatro planos y \fmm{M^*} la ampliada.
\end{fmprop}

\Needspace{7\baselineskip}
\begin{fmex}{Ejemplo resuelto}
\begin{enumerate}[label=\textbf{\arabic*.},leftmargin=*]
\item Estudia la posición relativa de \fmm{r:(x,y,z)=(1,0,2)+\lambda(1,-1,2)} y \fmm{s:(x,y,z)=(2,1,0)+\mu(2,1,-1)}.

\fmm{\vec d_r=(1,-1,2)} y \fmm{\vec d_s=(2,1,-1)} no son proporcionales: las rectas se cortan o se cruzan.

\fmm{\overrightarrow{PQ}=(1,1,-2)}, \fmm{\quad\begin{vmatrix}1&-1&2\\2&1&-1\\1&1&-2\end{vmatrix}=1\cdot(-2+1)+1\cdot(-4+1)+2\cdot(2-1)=-1-3+2=-2\neq0}

Los rangos son 2 y 3: las rectas \textbf{se cruzan}.
\item Estudia la posición relativa de \fmm{r:(x,y,z)=(1,0,-1)+\lambda(1,2,1)} y \fmm{s:(x,y,z)=(1,3,-2)+\mu(1,-1,2)} y, si se cortan, halla el punto de corte.

\fmm{\overrightarrow{PQ}=(0,3,-1)}, \fmm{\quad\begin{vmatrix}1&2&1\\1&-1&2\\0&3&-1\end{vmatrix}=1\cdot(1-6)-2\cdot(-1-0)+1\cdot(3-0)=-5+2+3=0}

Los directores no son proporcionales y el determinante es 0: son \textbf{secantes}.

\fmm{\begin{cases}1+\lambda=1+\mu\\2\lambda=3-\mu\\-1+\lambda=-2+2\mu\end{cases}\Rightarrow\lambda=\mu,\ 3\lambda=3\Rightarrow\lambda=\mu=1}. La tercera: \fmm{0=0\ \checkmark}. Punto de corte: \fmm{(2,2,0)}.
\end{enumerate}
\end{fmex}

\Needspace{16\baselineskip}
\subsection*{B. Recta y plano}

\Needspace{7\baselineskip}
\begin{fmprop}{Con el vector normal del plano y el director de la recta}
Sean \fmm{r} (punto \fmm{P}, director \fmm{\vec d}) y \fmm{\pi:Ax+By+Cz+D=0} (normal \fmm{\vec n}):

\begin{center}\renewcommand{\arraystretch}{2.5}\small\begin{tabularx}{\linewidth}{>{\raggedright\arraybackslash}p{3.2cm}>{\raggedright\arraybackslash}X>{\raggedright\arraybackslash}X}\toprule
\textbf{Condición} & \textbf{Posición} & \textbf{Rangos ($r$ en implícitas)}\\\midrule
$\vec n\cdot\vec d\neq0$ & \textbf{Secantes}: se cortan en un punto & $\mathrm{rg}\,M=\mathrm{rg}\,M^*=3$\\[3pt]
$\vec n\cdot\vec d=0,\ P\notin\pi$ & \textbf{Paralelos} & $\mathrm{rg}\,M=2,\ \mathrm{rg}\,M^*=3$\\[3pt]
$\vec n\cdot\vec d=0,\ P\in\pi$ & \textbf{La recta está contenida en el plano} & $\mathrm{rg}\,M=\mathrm{rg}\,M^*=2$\\[3pt]
\bottomrule\end{tabularx}\end{center}

Para hallar el punto de corte, sustituimos las paramétricas de la recta en la ecuación del plano y despejamos \fmm{\lambda}.
\end{fmprop}

\Needspace{16\baselineskip}
\subsection*{C. Dos planos}

\Needspace{7\baselineskip}
\begin{fmprop}{Comparando los coeficientes}
Sean \fmm{\pi_1:A_1x+B_1y+C_1z+D_1=0} y \fmm{\pi_2:A_2x+B_2y+C_2z+D_2=0}:

\begin{center}\renewcommand{\arraystretch}{2.5}\small\begin{tabularx}{\linewidth}{>{\raggedright\arraybackslash}p{5.2cm}>{\raggedright\arraybackslash}X>{\raggedright\arraybackslash}X}\toprule
\textbf{Condición} & \textbf{Rangos} & \textbf{Posición}\\\midrule
$\dfrac{A_1}{A_2}=\dfrac{B_1}{B_2}=\dfrac{C_1}{C_2}=\dfrac{D_1}{D_2}$ & $\mathrm{rg}\,M=\mathrm{rg}\,M^*=1$ & \textbf{Coincidentes}\\[3pt]
$\dfrac{A_1}{A_2}=\dfrac{B_1}{B_2}=\dfrac{C_1}{C_2}\neq\dfrac{D_1}{D_2}$ & $\mathrm{rg}\,M=1,\ \mathrm{rg}\,M^*=2$ & \textbf{Paralelos}\\[3pt]
$A_1,B_1,C_1$ no proporcionales a $A_2,B_2,C_2$ & $\mathrm{rg}\,M=\mathrm{rg}\,M^*=2$ & \textbf{Secantes}: se cortan en una recta\\[3pt]
\bottomrule\end{tabularx}\end{center}
\end{fmprop}

\Needspace{7\baselineskip}
\begin{fmex}{Ejemplo resuelto}
\begin{enumerate}[label=\textbf{\arabic*.},leftmargin=*]
\item Estudia la posición de \fmm{r:(x,y,z)=(1,0,2)+\lambda(2,1,-1)} y \fmm{\pi:x+y+z-6=0}.

\fmm{\vec n\cdot\vec d=(1,1,1)\cdot(2,1,-1)=2\neq0}. Son \textbf{secantes}.

\fmm{(1+2\lambda)+\lambda+(2-\lambda)-6=0\Rightarrow2\lambda-3=0\Rightarrow\lambda=\dfrac32}

Punto de corte: \fmm{\ \left(1+3,\ \dfrac32,\ 2-\dfrac32\right)=\left(4,\ \dfrac32,\ \dfrac12\right)}
\item Estudia la posición de los planos \fmm{\pi_1:x-2y+z=3}, \fmm{\pi_2:2x-4y+2z=1} y \fmm{\pi_3:-2x+4y-2z+6=0}.

\fmm{\pi_1,\pi_2}: \fmm{\ \dfrac12=\dfrac{-2}{-4}=\dfrac12\neq\dfrac31\Rightarrow} \textbf{paralelos}.

\fmm{\pi_3} se escribe \fmm{-2x+4y-2z=-6}, es decir, \fmm{-2\cdot(x-2y+z)=-2\cdot3}: \fmm{\pi_1} y \fmm{\pi_3} son \textbf{coincidentes}.

\fmuline{Cuidado}: para comparar el término independiente, ambos planos tienen que estar escritos de la misma forma (\fmm{D} en el mismo lado de la igualdad).
\end{enumerate}
\end{fmex}

\Needspace{18\baselineskip}
\section{Posición relativa de tres planos}\hypertarget{sec10}{}

\Needspace{7\baselineskip}
\begin{fmprop}{Teorema de Rouché aplicado a tres planos}
Los puntos comunes de tres planos son las soluciones del sistema que forman sus ecuaciones. Sean \fmm{M} la matriz de coeficientes y \fmm{M^*} la ampliada:

\fmdisp{\begin{cases}A_1x+B_1y+C_1z=D_1\\A_2x+B_2y+C_2z=D_2\\A_3x+B_3y+C_3z=D_3\end{cases}\qquad M=\begin{pmatrix}A_1&B_1&C_1\\A_2&B_2&C_2\\A_3&B_3&C_3\end{pmatrix}\qquad M^*=\left(\begin{array}{ccc|c}A_1&B_1&C_1&D_1\\A_2&B_2&C_2&D_2\\A_3&B_3&C_3&D_3\end{array}\right)}

\begin{center}\renewcommand{\arraystretch}{1.5}\small\begin{tabularx}{\linewidth}{>{\raggedright\arraybackslash}p{3.2cm}>{\raggedright\arraybackslash}X>{\raggedright\arraybackslash}X>{\raggedright\arraybackslash}X}\toprule
\textbf{$\mathrm{rg}\,M$} & \textbf{$\mathrm{rg}\,M^*$} & \textbf{Sistema} & \textbf{Posición de los planos}\\\midrule
3 & 3 & SCD & \textbf{Secantes en un punto}\\[3pt]
2 & 3 & SI & Ningún par de filas proporcional en $M$: \textbf{secantes dos a dos} (forman un prisma triangular)\\[3pt]
& & & Dos filas de $M$ proporcionales: \textbf{dos planos paralelos cortados por el tercero}\\[3pt]
2 & 2 & SCI (1 parámetro) & Ningún par de filas de $M^*$ proporcional: \textbf{secantes en una recta} (son de un mismo haz)\\[3pt]
& & & Dos filas de $M^*$ proporcionales: \textbf{dos planos coincidentes y el tercero secante}\\[3pt]
1 & 2 & SI & Ningún par de filas de $M^*$ proporcional: \textbf{tres planos paralelos}\\[3pt]
& & & Dos filas de $M^*$ proporcionales: \textbf{dos planos coincidentes y el tercero paralelo}\\[3pt]
1 & 1 & SCI (2 parámetros) & \textbf{Tres planos coincidentes}\\[3pt]
\bottomrule\end{tabularx}\end{center}

En el laboratorio de tres planos puedes ver y girar cada uno de los ocho casos.
\end{fmprop}

\Needspace{7\baselineskip}
\begin{fmex}{Ejemplo resuelto}
\begin{enumerate}[label=\textbf{\arabic*.},leftmargin=*]
\item Estudia la posición relativa de \fmm{\pi_1:x+y+z=1}, \fmm{\pi_2:x-y+2z=0} y \fmm{\pi_3:2x+3z=1}.

\fmm{|M|=\begin{vmatrix}1&1&1\\1&-1&2\\2&0&3\end{vmatrix}=1\cdot(-3)-1\cdot(3-4)+1\cdot(0+2)=-3+1+2=0}, y el menor \fmm{\begin{vmatrix}1&1\\1&-1\end{vmatrix}=-2\neq0}: \fmm{\mathrm{rg}\,M=2}.

Observamos que \fmm{\pi_3=\pi_1+\pi_2} (también el término independiente: \fmm{1+0=1}). Por tanto, \fmm{\mathrm{rg}\,M^*=2}.

Ningún par de planos es proporcional: los tres planos son \textbf{secantes en una recta}. Si cambiáramos \fmm{\pi_3} por \fmm{2x+3z=5}, tendríamos \fmm{\mathrm{rg}\,M^*=3} y formarían un \textbf{prisma}.
\item Discute, según los valores de \fmm{a}, la posición relativa de los planos \fmm{x+y+az=1}, \fmm{x+ay+z=1} y \fmm{ax+y+z=1}.

\fmm{|M|=\begin{vmatrix}1&1&a\\1&a&1\\a&1&1\end{vmatrix}=-(a-1)^2(a+2)\Rightarrow|M|=0\iff a=1\ o\ a=-2}

Si \fmm{a\neq1} y \fmm{a\neq-2}: \fmm{\mathrm{rg}\,M=\mathrm{rg}\,M^*=3}, \textbf{secantes en un punto}.

Si \fmm{a=1}: los tres planos son \fmm{x+y+z=1}: \textbf{coincidentes} (\fmm{\mathrm{rg}\,M=\mathrm{rg}\,M^*=1}).

Si \fmm{a=-2}: \fmm{\mathrm{rg}\,M=2} y el menor de \fmm{M^*} \fmm{\begin{vmatrix}1&1&1\\1&-2&1\\-2&1&1\end{vmatrix}=-9\neq0}, así que \fmm{\mathrm{rg}\,M^*=3}. Como no hay dos filas de \fmm{M} proporcionales, los planos son \textbf{secantes dos a dos} (prisma).
\end{enumerate}
\end{fmex}

\Needspace{18\baselineskip}
\section{Haces de planos}\hypertarget{sec11}{}

\Needspace{7\baselineskip}
\begin{fmdef}{Definiciones}
\begin{itemize}
\item El \textbf{haz de planos paralelos} a \fmm{\pi:Ax+By+Cz+D=0} es el conjunto de todos los planos que tienen el mismo vector normal:

\fmdisp{Ax+By+Cz+k=0,\qquad k\in\mathbb R}

\item El \textbf{haz de planos secantes} de arista \fmm{r} es el conjunto de todos los planos que contienen la recta \fmm{r}. Si \fmm{r} viene dada por dos planos, \fmm{r:\begin{cases}\pi_1:A_1x+B_1y+C_1z+D_1=0\\\pi_2:A_2x+B_2y+C_2z+D_2=0\end{cases}}, el haz es

\fmdisp{\alpha\,(A_1x+B_1y+C_1z+D_1)+\beta\,(A_2x+B_2y+C_2z+D_2)=0,\qquad(\alpha,\beta)\neq(0,0)}

\item Normalmente dividimos por \fmm{\alpha} y escribimos \fmm{\pi_1+\lambda\,\pi_2=0}. \fmuline{Cuidado}: de esta manera perdemos el plano \fmm{\pi_2} (que corresponde a \fmm{\alpha=0}). Si las condiciones del problema no dan ningún valor de \fmm{\lambda}, hay que comprobar si la solución es \fmm{\pi_2}.
\end{itemize}

Los haces son muy útiles para hallar el plano que \textbf{contiene una recta dada en implícitas} y cumple otra condición: pasar por un punto, ser paralelo a una recta, ser perpendicular a un plano\dots{} Imponemos la condición y calculamos \fmm{\lambda}.
\end{fmdef}

\Needspace{7\baselineskip}
\begin{fmex}{Ejemplo resuelto}
\begin{enumerate}[label=\textbf{\arabic*.},leftmargin=*]
\item Halla el plano que contiene la recta \fmm{r:\begin{cases}x+y-z-2=0\\2x-y+z-1=0\end{cases}} y pasa por \fmm{P(1,2,3)}.

Haz de planos de arista \fmm{r}: \fmm{\ (x+y-z-2)+\lambda\,(2x-y+z-1)=0}

Imponemos que pase por \fmm{P}: \fmm{\ (1+2-3-2)+\lambda\,(2-2+3-1)=0\Rightarrow-2+2\lambda=0\Rightarrow\lambda=1}

\fmm{(x+y-z-2)+(2x-y+z-1)=0\Rightarrow3x-3=0\Rightarrow\boxed{x=1}}
\item Halla el plano que contiene la misma recta \fmm{r} y es paralelo a la recta \fmm{s} de vector director \fmm{\vec d_s=(1,1,0)}.

Vector normal del plano del haz: \fmm{\ \vec n=(1+2\lambda,\ 1-\lambda,\ -1+\lambda)}

Paralelo a \fmm{s\iff\vec n\perp\vec d_s}: \fmm{\ (1+2\lambda)+(1-\lambda)=0\Rightarrow\lambda=-2}

\fmm{(x+y-z-2)-2(2x-y+z-1)=0\Rightarrow-3x+3y-3z=0\Rightarrow\boxed{x-y+z=0}}
\item Halla el plano paralelo a \fmm{2x-y+3z=1} que pasa por \fmm{P(1,1,1)}.

\fmm{2x-y+3z+k=0}, \fmm{\ 2-1+3+k=0\Rightarrow k=-4\Rightarrow\boxed{2x-y+3z-4=0}}
\end{enumerate}
\end{fmex}

\Needspace{16\baselineskip}
\subsection*{Recta que pasa por un punto y corta a otra}

\Needspace{7\baselineskip}
\begin{fmprop}{Cómo se plantea}
Por un punto \fmm{P} pasan infinitas rectas que cortan una recta \fmm{r} (todas están en el plano que contiene \fmm{P} y \fmm{r}). Para fijar una, el enunciado siempre da una \textbf{condición más}. Hay dos métodos:

\begin{itemize}
\item \textbf{Punto genérico}: el punto de corte es \fmm{Q=A+\lambda\,\vec d_r\in r}. La recta buscada es \fmm{PQ}, de director \fmm{\overrightarrow{PQ}} (que depende de \fmm{\lambda}). Imponemos la condición a \fmm{\overrightarrow{PQ}} y hallamos \fmm{\lambda}:

\begin{itemize}
\item paralela a un plano \fmm{\pi}: \fmm{\ \overrightarrow{PQ}\cdot\vec n_\pi=0};
\item perpendicular a un vector \fmm{\vec u}: \fmm{\ \overrightarrow{PQ}\cdot\vec u=0} (si \fmm{\vec u=\vec d_r}, es la recta que corta a \fmm{r} \textbf{perpendicularmente}: lo veremos en la unidad 4);
\item que también corte otra recta \fmm{s}: \fmm{\ \overrightarrow{PQ}}, \fmm{\vec d_s} y un vector de \fmm{P} a un punto de \fmm{s} tienen que ser coplanarios.
\end{itemize}
\item \textbf{Planos}: la recta buscada está en el plano \fmm{\pi_r} que contiene \fmm{P} y \fmm{r} (\fmm{\vec n=\vec d_r\times\overrightarrow{AP}}, o un plano del haz de arista \fmm{r}). Si también tiene que cortar una recta \fmm{s}, está en el plano \fmm{\pi_s} que contiene \fmm{P} y \fmm{s}, y es \fmm{\pi_r\cap\pi_s}. O, más rápido: cortamos \fmm{s} con \fmm{\pi_r}; el punto de corte \fmm{Q} da la recta \fmm{PQ}.
\end{itemize}
\vspace*{0.2cm}
\begin{fmnote}\fmuline{Cuidado}: al final hay que comprobar que la recta \fmm{PQ} no es paralela a \fmm{r} (o a \fmm{s}); si lo fuera, no las cortaría y el problema no tendría solución.\end{fmnote}
\end{fmprop}

\Needspace{7\baselineskip}
\begin{fmex}{Ejemplo resuelto}
\begin{enumerate}[label=\textbf{\arabic*.},leftmargin=*]
\item Halla la recta que pasa por \fmm{P(1,0,2)}, corta la recta \fmm{r:(x,y,z)=(0,1,1)+\lambda(1,1,0)} y es paralela al plano \fmm{\pi:2x-y+z=0}.

Punto genérico de \fmm{r}: \fmm{\ Q(\lambda,\,1+\lambda,\,1)}, \fmm{\ \overrightarrow{PQ}=(\lambda-1,\ 1+\lambda,\ -1)}.

Paralela a \fmm{\pi}: \fmm{\ \overrightarrow{PQ}\cdot(2,-1,1)=2\lambda-2-1-\lambda-1=\lambda-4=0\Rightarrow\lambda=4}, \fmm{\ Q(4,5,1)}.

\fmm{\overrightarrow{PQ}=(3,5,-1)}, que no es paralelo a \fmm{\vec d_r}: \fmm{\ \boxed{(x,y,z)=(1,0,2)+t\,(3,5,-1)}}
\item Halla la recta que pasa por el origen y corta las rectas \fmm{r:(x,y,z)=(1,0,1)+\lambda(1,1,0)} y \fmm{s:(x,y,z)=(0,1,-1)+\mu(1,0,2)}.

Plano que contiene \fmm{O} y \fmm{r}: \fmm{\ \vec n=\vec d_r\times\overrightarrow{OA}=(1,1,0)\times(1,0,1)=(1,-1,-1)}, \fmm{\ \pi_r:x-y-z=0}.

Cortamos \fmm{s} con \fmm{\pi_r}: \fmm{\ \mu-(1)-(-1+2\mu)=-\mu=0\Rightarrow\mu=0}, \fmm{\ Q(0,1,-1)}.

Recta \fmm{OQ}: \fmm{\ \boxed{(x,y,z)=t\,(0,1,-1)}}. Comprobación: no es paralela a \fmm{r}, y el punto \fmm{(0,-1,1)} (\fmm{t=-1}) es de \fmm{r} (\fmm{\lambda=-1}) .
\end{enumerate}
\end{fmex}

  
\end{document}
